\chapter{분포 확인과 데이터 변환: Q-Q 플롯에서 정규화까지}
\label{ch:qq-normalization}
% source: deck 7 slide 1

\begin{inbrief}
  \item[Why] 모수를 추정하거나 정규성을 가정하는 검정을 쓰려면 데이터가 어떤
             분포를 따르는지 먼저 확인해야 하고, 단위와 크기가 서로 다른
             변수는 그대로 비교하거나 합칠 수 없다.
  \item[When] 새 데이터를 받아 분포 모형을 고르기 전, 거리 기반 알고리즘에
              넣거나 여러 실험의 측정값을 합치기 전, 결측값(Missing
              Value)이 보일 때.
  \item[Where] 영화 평점(\Cref{ch:hw2}), 유전자 발현(Gene Expression), 붓꽃
               꽃받침 길이(Sepal Length), 나이와 소득, 패스트푸드 당류
               (\texttt{sugars})와 열량(\texttt{calories}).
  \item[How] 히스토그램 $\to$ CDF $\to$ Q-Q 플롯(Quantile-Quantile Plot)
             순서로 후보 분포와 대 보고, 목적에 맞춰 이산화(Discretization),
             범위 정규화(Range Normalization), z-정규화(z-Normalization),
             분위수 정규화(Quantile Normalization), 로버스트 스케일링
             (Robust Scaling), 로그 변환(Log Transformation) 중 하나를 고른
             뒤, 결측값은 삭제하거나 대체(Imputation)한다.
\end{inbrief}

이 장의 출처는 \texttt{notes\_text/7\_qqPlotandNormalization/slides.md}
(26 슬라이드) 하나다. 슬라이드는 문헌을 밝히지 않으므로 새 참고문헌 항목을
만들지 않는다. 예외 하나(\Cref{sec:qq-future})는 그림 안에 인쇄된 출처를
그대로 옮긴다.

\analogy{지도 비유}를 이어 쓴다. 분포를 가정하는 일은 어떤 투영법을 쓸지
정하는 일과 같다. Q-Q 플롯은 그 투영법이 맞는지 격자선 몇 개를 실제 지형과
대 보는 일이고, 정규화는 축척이 다른 지도 여러 장을 같은 축척으로 맞추는
일이다.

% =====================================================================
\section{배경}
\label{sec:qq-background}
% source: deck 7 slide 2

이 장은 다음 개념을 그대로 가져와 쓴다.

\begin{itemize}
  \item 분위수와 선형 보간(\Cref{sec:quantile}, \Cref{def:quantile},
        \Cref{ex:quantile}): $n$개 값을 정렬한 뒤 위치 $(i-1)/(n-1)$에서
        보간한다. 이 장에는 \textbf{다른 두 규칙}이 더 나온다
        (\Cref{sec:qq}의 함정, \Cref{sec:discretization}의 함정).
  \item 누적분포함수(Cumulative Distribution Function, CDF). 정식 정의는
        \Cref{def:ecdf}(\Cref{ch:compare-dist})에서 다루지만, \Cref{sec:cdf}가
        아직 초안(stub)이므로 이 장이 두 문장으로 직접 요약한다. 값 $x$에서의
        CDF $F(x) = P(X \le x)$는 ``$x$ 이하가 전체에서 차지하는 비율''이다.
        데이터로 만든 계단 모양은 \term{경험적 CDF}(Empirical CDF,
        $\hat F$)라 부른다(정식 정의는 \Cref{def:ecdf}).
  \item 정규분포와 그 밀도(\Cref{sec:normal}, \Cref{def:normal}). 표준정규
        CDF는 $\Phi$, 그 역함수는 $\Phi^{-1}$로 쓴다.
  \item 네 후보 분포(멱법칙, 지수, 균등, 정규)의 모수 추정
        (\Cref{sec:hw2-estimators}, \Cref{tab:hw2-params}). 슬라이드 2의
        추정식 넷은 이 표의 식과 같으므로 다시 유도하지 않는다.
  \item 평균, 표준편차($n-1$로 나눔), 사분위 범위
        (\Cref{sec:center}, \Cref{def:spread}, \Cref{subsec:fence}). 이 장의
        슬라이드 18은 $\hat\sigma$를 $n$으로 나눈다(뒤의 함정 참고).
  \item 왜도와 그 부호(\Cref{sec:skewness}, \Cref{sec:skew-sign}). Q-Q 모양
        읽기와 로그 변환에 쓴다.
  \item 결측값과 ``0은 결측이 아니다''(\Cref{sec:missing},
        \Cref{ex:zero-as-missing}). \Cref{sec:impute}와 $\log 0$ 함정의
        선행 조건이다.
  \item 히스토그램 구간 나누기(\Cref{sec:histogram}, \Cref{subsec:binning}).
        이산화는 같은 연산을 \textbf{그리기}가 아니라 \textbf{변수
        바꾸기}에 쓰는 것이다.
\end{itemize}

새 개념 둘은 이 장에서 직접 정의한다. 역CDF $\hat F^{-1}(q)$는 비율 $q$를
넣으면 그 비율에 \textbf{처음 도달하는} 값을 돌려준다. 분위수를 CDF
방향에서 본 것과 같다. \term{플로팅 위치}(Plotting Position) $p_i$는 정렬된
$i$번째 값에 붙이는 ``이론적 위치''로, 이 장에서는 $p_i = i/(n+1)$을 쓴다
(\Cref{sec:qq}의 함정에서 \Cref{ex:quantile}의 규칙과 비교한다).

\Cref{tab:qq-symbols}가 이 장에서 쓰는 기호를 모은다.

\begin{table}[ht]\centering
\begin{tabularx}{\linewidth}{@{}l X@{}}
\toprule
기호 & 뜻 \\
\midrule
$x_i$, $x_{(i)}$ & 원래 값, 정렬 후 $i$번째로 작은 값 \\
$n$ & 표본 크기 \\
$p_i = i/(n+1)$ & 플로팅 위치(이론 분포와 비교할 때) \\
$\Phi$, $\Phi^{-1}$ & 표준정규 CDF, 그 역함수 \\
$\hat F$, $\hat F^{-1}(q)$ & 경험적 CDF, 그 역함수(비율 $q$의 분위수) \\
$k$, $w$, $v_i$ & 이산화 구간 수, 구간 폭, 구간 경계 \\
$\hat r$ & 표본 범위(최댓값 $-$ 최솟값) \\
$\hat\mu$, $\hat\sigma$ & 표본 평균, 표본 표준편차 \\
$x_i'$ & 변환 후 값 \\
$\IQR$ & 사분위 범위 \\
\bottomrule
\end{tabularx}
\caption{이 장에서 쓰는 기호.}
\label{tab:qq-symbols}
\end{table}

\begin{pitfall}[기호 충돌 둘]
슬라이드 2는 정규분포의 분산을 $\lambda$로 적어 지수분포의
모수 $\lambda$(\Cref{sec:exponential})와 겹친다. 이 책은 $\sigma^2$(또는
$\hat\sigma^2$)로 쓴다. 같은 슬라이드의 균등분포 그림은 구간을
$[\alpha, \beta]$로 적어 멱법칙 지수 $\alpha$(\Cref{sec:hw2-estimators})와
겹친다. 이 책은 균등분포 구간을 같은 슬라이드의 다른 그림처럼 $[a, b]$로 통일해
쓴다.
\end{pitfall}

% =====================================================================
\section{왜 필요한가}
\label{sec:qq-why}
% source: deck 7 slide 2, 12

\begin{itemize}
  \item 문제 1(슬라이드 2): 모수를 추정하려면 먼저 분포의 \textbf{모양}을
        정해야 한다. 같은 영화 평점에 네 후보(멱법칙, 지수, 균등, 정규)의
        모수를 모두 구할 수 있지만(\Cref{tab:hw2-params}), 모수를 구했다고
        모양이 맞는 것은 아니다.
  \item 문제 2: 뒤 장의 t검정, F검정은 정규성을 가정한다(\Cref{sec:cd-t-one},
        \Cref{sec:cd-future}). 가정을 확인할 도구가 없으면 검정 결과를 믿을
        근거가 없다.
  \item 문제 3(슬라이드 12, 16): 분포를 확인한 다음에도 변수의 단위, 크기,
        모양이 분석에 맞지 않을 수 있다. 소득이 나이를 압도하고(슬라이드
        18), 긴 꼬리가 평균을 끌고 간다(\Cref{sec:skew-sign}).
\end{itemize}

처방은 셋으로 나뉜다. 앞 절반은 \textbf{확인}(히스토그램 $\to$ CDF $\to$
Q-Q), 뒤 절반은 \textbf{변환}(이산화, 정규화, 로그), 마지막은
\textbf{빈칸 처리}(결측값)다. 슬라이드 12의 변환 목록(이산화, 정규화,
중심화와 스케일링, 로그)이 뒤 절반의 목차다.

% =====================================================================
\section{히스토그램에 후보 분포 겹치기}
\label{sec:dist-check-hist}
% source: deck 7 slide 3

\begin{inbrief}
  \item[Why] 모양이 전혀 다른 후보는 한눈에 걸러진다.
  \item[When] 후보가 여럿이고 모수를 이미 구했을 때.
  \item[Where] 영화 평점 네 후보, 패스트푸드 열량의 정규 적합.
  \item[How] 밀도(Density) 척도 히스토그램 위에, 추정한 모수로 만든 밀도
             곡선을 겹쳐 그린다.
\end{inbrief}

\paragraph{고치는 문제.} \Cref{sec:hw2-estimators}처럼 모수를 최대우도로
구해도, 그 모수가 만드는 곡선이 데이터의 \textbf{모양}과 맞는지는 따로
봐야 한다. 모수는 항상 계산되지만, 곡선과 히스토그램의 일치 여부는
계산으로 보장되지 않는다.

\paragraph{아이디어.} 데이터를 밀도 척도(넓이 합 1)로 그린 히스토그램 위에,
각 후보 분포의 밀도 곡선을 같은 좌표에 겹친다. 곡선이 막대의 윤곽을 따라가면
그 후보가 유력하고, 막대와 곡선의 모양 자체가 다르면(단조감소 대 종 모양처럼)
모수를 아무리 조정해도 맞지 않는다.

\Cref{fig:qq-four-pdf}가 영화 평점(\Cref{sec:hw2-movie})에 네 후보를 겹친다.

\begin{figure}[ht]\centering
  \includegraphics[width=0.95\linewidth]{07_01_movie_four_pdfs.png}
  \caption{영화 평점(4392편) 히스토그램(0.1 격자) 위에 네 후보 밀도.
  멱법칙 $\hat\alpha = 1.8505$($x_{\min}=1.9$), 지수 $\hat\lambda = 0.1606$,
  균등 $[1.9, 8.5]$, 정규 $\hat\mu = 6.227$, $\hat\sigma = 0.893$
  (\Cref{tab:hw2-params}와 일치).}
  \label{fig:qq-four-pdf}
\end{figure}

\begin{pitfall}[모수를 구했다고 모양이 맞는 것은 아니다]
영화 평점에 지수분포를 최대우도로 맞추면 $\hat\lambda = 0.1606$은 올바른
추정값이지만, 지수분포의 밀도는 $x=0$ 근처에서 가장 높고 단조감소한다.
영화 평점은 6.3 근처에서 봉우리를 이루는 종 모양이라, 지수분포는 이
네 후보 중 최악이다(\Cref{tab:hw2-movie-verdict}, $D = 0.490$). 모수
추정이 맞았다는 사실이 모형 선택이 맞았다는 뜻은 아니다.
\end{pitfall}

\Cref{fig:normal-fit}이 패스트푸드 \texttt{calories}에 정규 곡선을 겹친다.
이 그림의 $\hat\mu, \hat\sigma$는 \texttt{scipy.stats.norm.fit}이 주는
최대우도 값($n$으로 나눔)이다.

\begin{figure}[ht]\centering
  \includegraphics[width=0.7\linewidth]{06_01_normal_fit.png}
  \caption{패스트푸드 \texttt{calories} 히스토그램과 최대우도(MLE) 정규
  곡선. 히스토그램은 종 모양에 가까워 보이지만, \Cref{sec:qq-apply}의
  Q-Q 플롯과 Shapiro-Wilk 검정은 정규성을 기각한다.}
  \label{fig:normal-fit}
\end{figure}

% =====================================================================
\section{CDF로 한 번 더 확인}
\label{sec:dist-check-cdf}
% source: deck 7 slide 4

\begin{inbrief}
  \item[Why] 히스토그램은 구간 폭에 따라 모양이 바뀌지만, CDF는 구간이
             필요 없다.
  \item[When] 히스토그램으로 후보를 한두 개로 줄인 뒤.
  \item[Where] 영화 평점 네 후보의 CDF.
  \item[How] 경험적 CDF 계단과 후보 CDF를 겹쳐, 가장 크게 벌어지는 지점을
             본다.
\end{inbrief}

\paragraph{고치는 문제.} 히스토그램은 구간 경계를 어디에 놓느냐에 따라
막대의 높이가 흔들린다(\Cref{subsec:binning}). 두 후보의 밀도가 비슷하면
구간 하나 차이로 판정이 뒤집힐 수 있다.

\paragraph{아이디어.} 구간을 나누지 않는 CDF로 같은 비교를 한다. 경험적
CDF $\hat F$와 후보의 이론 CDF를 같은 축에 겹치면, 가장 크게 벌어지는
세로 거리가 두 곡선이 가장 다른 지점이다. 그 최대 세로 격차가
\term{콜모고로프-스미르노프 통계량}(KS Statistic) $D$다(\Cref{def:ks-statistic}).

\begin{figure}[ht]\centering
  \includegraphics[width=0.95\linewidth]{07_02_movie_four_cdfs.png}
  \caption{영화 평점의 경험적 CDF 계단과 네 후보 CDF. 멱법칙, 지수
  패널은 로그 가로축(슬라이드 4와 같음). 패널마다 한 표본에서 잰 최대
  세로 격차 위치에 세로선을 표시했다. 이 격차는 한 표본에서 잰 값이라,
  \Cref{tab:hw2-movie-verdict}의 두 표본 KS 거리(모형 표본과 비교한
  $D$)와는 다를 수 있다.}
  \label{fig:qq-four-cdf}
\end{figure}

\begin{pitfall}[같은 데이터로 추정한 모수의 낙관 편향]
이 장의 후보는 데이터 \textbf{같은 것}으로 모수를 추정했다. 그 상태로
KS p값을 계산하면 실제보다 낙관적으로(기각이 잘 안 되는 쪽으로) 나온다
(\Cref{sec:normal}의 함정). 이 절은 CDF를 \textbf{눈으로 보는} 도구로만
쓰고, p값을 계산하지 않는다.
\end{pitfall}

% =====================================================================
\section{Q-Q 플롯}
\label{sec:qq}
% source: deck 7 slide 5, 6, 7, 8, 9, 10

\begin{inbrief}
  \item[Why] CDF는 꼬리(0과 1 근처)가 눌려 보여 꼬리 차이를 놓친다. Q-Q는
             꼬리를 펼쳐 보인다.
  \item[When] 후보 하나가 남았고 ``어디서 어긋나는가''를 알고 싶을 때.
  \item[Where] 점 15개의 정규 대 균등 비교, 치우침과 꼬리의 네 모양.
  \item[How] 데이터를 정렬 $\to$ 플로팅 위치 $p_i = i/(n+1)$에서 후보의
             분위수 $\Phi^{-1}(p_i)$를 구함 $\to$ 점
             $(\Phi^{-1}(p_i), x_{(i)})$를 찍고 직선에서 벗어난 모양을
             읽는다.
\end{inbrief}

\paragraph{고치는 문제.} CDF는 $[0,1]$ 사이의 값이므로, 확률이 아주 작거나
큰 영역(꼬리)에서는 세로축이 눌려 보인다. 꼬리가 두껍거나 얇은 차이는
CDF 그림에서 거의 티가 나지 않는다.

\paragraph{아이디어.} 세로축을 확률 대신 \textbf{값 자체}로 바꾼다. 두
분포(데이터, 후보)에서 같은 위치의 분위수를 짝지어 점으로 찍으면, 두 분포가
같을 때 그 점들은 직선 $y = \hat\mu + \hat\sigma z$ 위에 놓인다. 직선에서
벗어난 방향과 크기가 ``어디서, 어떻게'' 다른지 말해준다.

\Cref{sec:hw2-compare-how}가 이미 같은 도구(\term{분위수 대 분위수
도표}(Quantile-Quantile Plot, QQ Plot))를 정의했다. 그 장은 가로축을
데이터, 세로축을 모형 표본으로 뒀지만, 이 장의 슬라이드는 반대로
가로축을 이론 분위수, 세로축을 데이터로 둔다(아래 함정 (2)).

\begin{definition}[Q-Q 플롯]
\label{def:qq-plot}
정렬된 데이터 $x_{(1)} \le \dots \le x_{(n)}$과 플로팅 위치
$p_i = i/(n+1)$에 대해, 점
\[
  \bigl(\Phi^{-1}(p_i),\ x_{(i)}\bigr), \qquad i = 1, \dots, n
\]
을 찍은 산점도를 Q-Q 플롯이라 한다. 후보 분포가 정규분포가 \textbf{아닐} 때는
$\Phi^{-1}$을 그 분포의 역CDF로 바꿔 쓴다(예: 균등 $[0,1]$이면
$\Phi^{-1}(p_i) = p_i$).
데이터가 그 후보 분포에서 왔다면, 점은 한 직선 위에 놓인다. 정규분포 후보라면
그 직선이 $y = \hat\mu + \hat\sigma z$이고, 다른 후보라면 그 분포의 위치와
척도 모수로 만든 직선이다.
\end{definition}

\paragraph{손 계산: 점 5개.} 데이터 $\{3, 4, 5, 6, 12\}$(이미 정렬됨,
책의 예시 값)로 정규 Q-Q를 직접 만든다.

\begin{align*}
  &p_i = i/6,\ i=1..5 && \text{플로팅 위치}\\
  &(0.1667,\ 0.3333,\ 0.5,\ 0.6667,\ 0.8333) && \text{다섯 값}\\
  &\Phi^{-1}(p_i) = (-0.9674,\ -0.4307,\ 0,\ 0.4307,\ 0.9674) && \text{\texttt{norm.ppf}}\\
  &\bar x = (3{+}4{+}5{+}6{+}12)/5 = 6 && \text{평균}\\
  &(3{-}6)^2+(4{-}6)^2+(5{-}6)^2+(6{-}6)^2+(12{-}6)^2 = 50 && \text{제곱편차 합}\\
  &s = \sqrt{50/4} = 3.536 && \text{표준편차($n-1$)}
\end{align*}

기준선은 $y = 6 + 3.536\,z$다. \Cref{tab:qq-steps}가 각 점의 위치와 기준선
위의 값을 정리한다.

\begin{table}[ht]\centering
\begin{tabularx}{\linewidth}{@{}l X X X X@{}}
\toprule
순위 & 값 $x_{(i)}$ & $p_i$ & $\Phi^{-1}(p_i)$ & 기준선 값 $6+3.536z$ \\
\midrule
1 & 3  & 0.1667 & $-0.9674$ & 2.58 \\
2 & 4  & 0.3333 & $-0.4307$ & 4.48 \\
3 & 5  & 0.5    & 0         & 6.00 \\
4 & 6  & 0.6667 & 0.4307    & 7.52 \\
5 & 12 & 0.8333 & 0.9674    & 9.42 \\
\bottomrule
\end{tabularx}
\caption{점 5개의 Q-Q 손 계산(책의 예시 값).}
\label{tab:qq-steps}
\end{table}

값 3, 12는 기준선보다 위(데이터가 기준선 값 2.58, 9.42보다 크다), 값 4, 6은
기준선보다 아래다. 양 끝이 위, 가운데가 아래인 모양은 \textbf{아래로 볼록}이고,
이는 \Cref{tab:qq-shapes}에서 오른쪽 치우침(Right Skew)을 뜻한다. 값 12
하나가 이 모양을 만든다.

\Cref{fig:qq-steps}가 같은 절차를 정규와 균등 두 경우로 그린다.

\begin{figure}[ht]\centering
  \includegraphics[width=0.95\linewidth]{07_03_qq_steps.png}
  \caption{Q-Q 플롯 세 단계, 정규(위 줄)와 균등 $[0,1]$(아래 줄, 이론
  분위수는 $p_i$ 그대로). (1) 데이터 점 (2) 이론 분위수 선(가운데가
  좁고 가장자리가 넓음, 간격은 가운데 $0.431$, 가장자리 $0.537$)
  (3) Q-Q 점과 기준선. 데이터 $\{3,4,5,6,12\}$(책의 예시 값).}
  \label{fig:qq-steps}
\end{figure}

\paragraph{가운데가 좁고 가장자리가 넓은 이유.} 슬라이드 6은 표본 15개
예제에서 이 사실을 그림으로만 보인다. 표준정규 분위수로 직접 확인하면

\begin{align*}
  &\Phi^{-1}(9/16) - \Phi^{-1}(8/16) = 0.1573 - 0 = 0.1573 && \text{가운데 간격}\\
  &\Phi^{-1}(2/16) - \Phi^{-1}(1/16) = -1.1503 - (-1.5341) = 0.3838 && \text{가장자리 간격}
\end{align*}

가장자리 간격이 가운데 간격의 약 2.4배다. 정규분포의 밀도가 가운데에서
높고 꼬리에서 낮으므로, 같은 확률 폭 $1/16$을 담는 데 필요한 $x$축 폭이
가장자리에서 더 넓다.

\begin{pitfall}[플로팅 위치: 소프트웨어마다 규칙이 다르다]
이 장은 $p_i = i/(n+1)$을 쓴다(값을 보면 규칙이 그렇지만, 슬라이드
자체는 규칙을 밝히지 않는다). \Cref{ex:quantile}의 규칙은
$(i-1)/(n-1)$이다. \texttt{scipy.stats.probplot}은 이 둘 중 어느 것도
아닌 \textbf{Filliben의 근사식}을 쓴다.
\[
  p_1 = 1 - 0.5^{1/n}, \qquad p_n = 0.5^{1/n}, \qquad
  p_i = \frac{i - 0.3}{n + 0.4}\ (1 < i < n).
\]
$n=15$에서 이 장의 규칙은 $p_1 = i/(n+1) = 0.0625$, Filliben은
$p_1 = 1-0.5^{1/15} = 0.0452$로, 직접 확인하면
(\texttt{scipy.stats.probplot(x, dist="norm")}가 반환하는 이론 분위수를
$\Phi$로 되돌려 계산) 값이 다르다. $\Phi^{-1}$ 값도 $-1.534$ 대 $-1.694$로
벌어진다. 어느 규칙을 쓰든 Q-Q의 \textbf{모양} 판단(직선인가 아닌가, 어느
쪽으로 휘는가)은 바뀌지 않는다.
\end{pitfall}

\begin{pitfall}[축의 방향이 장마다 다르다]
이 장의 슬라이드는 가로축 = 이론 분위수, 세로축 = 데이터다.
\Cref{sec:hw2-compare-how}(\Cref{ch:hw2})는 가로축 = 데이터, 세로축 =
모형 표본이다. 두 장의 정의(\Cref{def:qq-plot})는 같지만 축을 바꿔
그렸으므로, 모양을 읽는 방향(위로 볼록/아래로 볼록)이 서로 뒤집힌다.
Q-Q 플롯을 볼 때는 항상 축 이름부터 읽는다.
\end{pitfall}

\begin{pitfall}[작은 표본으로 가른 판정]
슬라이드 8은 점 15개짜리 Q-Q 둘(정규, 균등)을 나란히 놓고 ``균등이 더
잘 맞는다''고 결론 낸다. 점이 15개뿐이면 두 후보를 가를 통계적 힘이
약하다. 그림 하나로 내린 판정은 $n$과 함께 말해야 한다.
\end{pitfall}

\paragraph{네 가지 모양.} \Cref{tab:qq-shapes}가 Q-Q 점의 휘는 방향과
히스토그램 모양을 짝짓는다.

\begin{table}[ht]\centering
\begin{tabularx}{\linewidth}{@{}l X X@{}}
\toprule
모양 & Q-Q 점(가로 이론, 세로 데이터) & 이 책의 다른 자리 \\
\midrule
정규 & 직선 & 없음 \\
오른쪽 치우침(Right Skew) & 아래로 볼록(양 끝이 선 위로 휨) & \Cref{fig:skew-right} \\
왼쪽 치우침(Left Skew) & 위로 볼록(양 끝이 선 아래로) & \Cref{fig:skew-left} \\
두꺼운 꼬리(Fat Tails) & 왼쪽 끝은 선 아래, 오른쪽 끝은 선 위(S자) & 없음 \\
얇은 꼬리(Thin Tails) & 왼쪽 끝은 선 위, 오른쪽 끝은 선 아래(뒤집힌 S자) & 없음 \\
\bottomrule
\end{tabularx}
\caption{Q-Q 점의 모양 읽기.}
\label{tab:qq-shapes}
\end{table}

\begin{figure}[ht]\centering
  \includegraphics[width=0.95\linewidth]{07_04_qq_shapes.png}
  \caption{다섯 모양의 히스토그램과 정규 Q-Q(각 $n=1000$, 시드 6330).
  위에서부터 정규, 오른쪽 치우침(로그정규 표본), 왼쪽 치우침(부호를
  바꾼 표본), 두꺼운 꼬리($t$분포, 자유도 3), 얇은 꼬리(균등분포).
  히스토그램 위 정규 곡선은 각 표본의 평균과 표준편차로 맞췄다.}
  \label{fig:qq-shapes}
\end{figure}

\begin{pitfall}[얇은 꼬리 예시가 약하면 뒤집힌 S자가 잘 안 보인다]
슬라이드 10의 얇은 꼬리 Q-Q는 거의 직선처럼 보이고,
겹쳐 그린 정규 곡선도 데이터의 표준편차에 맞춰져 있지 않다.
\Cref{fig:qq-shapes}는 얇은 꼬리 표본으로 균등분포를 써서 뒤집힌 S자
모양이 보이게 그렸다.
\end{pitfall}

% =====================================================================
\section{두 데이터셋의 Q-Q}
\label{sec:qq-two-samples}
% source: deck 7 slide 5, 11

\begin{inbrief}
  \item[Why] 비교 대상이 이론 분포가 아니라 다른 표본일 때도 모양을
             비교해야 한다.
  \item[When] 두 표본의 분포가 같은지 볼 때, 표본 크기가 달라도.
  \item[Where] 원래 데이터 15개 대 새 데이터 4개.
  \item[How] 작은 쪽 표본의 위치마다 큰 쪽 표본의 분위수를 보간으로 구해
             짝짓는다.
\end{inbrief}

\paragraph{고치는 문제.} \Cref{sec:qq}는 후보가 \textbf{이론 분포}일 때만
다뤘다. 그러나 비교 대상이 다른 \textbf{표본}이고, 그 표본의 크기가 원래
데이터와 다르면(슬라이드 5), 위치를 어떻게 짝지을지가 새 문제다.

\paragraph{아이디어.} 두 표본의 크기가 다르면, 작은 쪽 표본의 각 위치에
대응하는 값을 큰 쪽 표본에서 \textbf{보간}으로 만든다. 그러면 표본 크기가
달라도 같은 개수의 점을 찍을 수 있다.

\begin{example}[크기가 다른 두 표본]
\label{ex:qq-two-samples}
원래 표본 A($n=15$, 정렬): 2, 3, 3.5, 4, 4.5, 5, 5, 5.5, 6, 6, 6.5, 7, 8,
9, 11. 새 표본 B($n=4$): 4, 5.5, 7, 9.5(책의 예시 값). B의 위치를
\Cref{ex:quantile}의 규칙 $(i-1)/(n-1)$로 잡으면 $0, 1/3, 2/3, 1$이다.

\begin{align*}
  &p = 0 \to \text{A의 인덱스 } 0 \to 2 && \text{최솟값끼리}\\
  &p = 1/3 \to \text{인덱스 } 4.667 \to 4.5 + 0.667(5{-}4.5) = 4.833 && \text{보간}\\
  &p = 2/3 \to \text{인덱스 } 9.333 \to 6 + 0.333(6.5{-}6) = 6.167 && \text{보간}\\
  &p = 1 \to \text{인덱스 } 14 \to 11 && \text{최댓값끼리}
\end{align*}

점 (B, A): $(4, 2)$, $(5.5, 4.833)$, $(7, 6.167)$, $(9.5, 11)$. 기울기가
1보다 크므로 A가 B보다 더 퍼져 있다.
\end{example}

\begin{keyidea}[규칙을 일부러 바꾼 이유]
\Cref{sec:qq}의 손 계산은 이론 분포와 비교하므로 $p_i = i/(n+1)$을
썼다. 정규 분위수는 $p=0, 1$에서 무한대라 끝점을 쓸 수 없기 때문이다.
\Cref{ex:qq-two-samples}는 두 표본끼리 비교하므로 \Cref{ex:quantile}의
$(i-1)/(n-1)$을 써서 B의
최솟값, 최댓값을 A의 최솟값, 최댓값과 짝짓는다. 같은 B에 $i/(n+1)$을
쓰면 A의 $0.2, 0.4, 0.6, 0.8$ 분위수(값 $3.9, 5.0, 6.0, 7.2$)와 짝지어져
다른 그림이 나온다. 두 규칙 모두 정당하고, \textbf{어느 것을 썼는지
밝히는 것}이 요점이다.
\end{keyidea}

\begin{figure}[ht]\centering
  \includegraphics[width=0.95\linewidth]{07_05_qq_two_samples.png}
  \caption{왼쪽: A(15개) 세로줄에 B(4개)의 보간 위치를 점선으로 표시.
  오른쪽: 점 $(B, A)$ 네 개와 기준선 $y=x$. \Cref{ex:qq-two-samples}의
  값.}
  \label{fig:qq-two-samples}
\end{figure}

% =====================================================================
\section{이산화}
\label{sec:discretization}
% source: deck 7 slide 12, 13, 14, 15

\begin{inbrief}
  \item[Why] 수치형 속성을 범주로 바꿔야 쓸 수 있는 방법(빈도표, 범주형
             모형)이 있다.
  \item[When] 연속값을 몇 개 등급으로 묶을 때.
  \item[Where] 붓꽃 꽃받침 길이를 4구간으로.
  \item[How] 등간격(Equal-Width) $v_i = x_{\min} + i w$,
             $w = (x_{\max}-x_{\min})/k$, 또는 등빈도(Equal-Frequency)
             $v_i = \hat F^{-1}(i/k)$.
\end{inbrief}

\paragraph{고치는 문제.} 히스토그램(\Cref{sec:histogram})은 값을 구간으로
묶어 \textbf{그린다}. 그런데 빈도표나 범주형 모형처럼 구간으로 묶은 결과
자체를 \textbf{새 변수}로 써야 할 때가 있다. 그리기와 변수 바꾸기는
같은 연산(구간 나누기)을 다른 목적에 쓰는 것이다.

\paragraph{아이디어.} 구간을 나누는 규칙은 두 가지다. 폭을 똑같이 맞추는
등간격, 또는 개수를 똑같이 맞추는 등빈도. 등간격은 $x$축 위에서 폭
$w=(x_{\max}-x_{\min})/k$인 $k$개 구간을 만들고, 등빈도는 경험적 역CDF
$\hat F^{-1}(i/k)$로 개수가 같아지는 경계를 찾는다.

\begin{example}[붓꽃 꽃받침 길이 4구간]
\label{ex:qq-iris}
붓꽃 데이터(\Cref{sec:qq-iris-source}) $n=150$, 범위 $[4.3, 7.9]$다.

등빈도(슬라이드 15): $\hat F^{-1}(0.25) = 5.1$, $\hat F^{-1}(0.5) = 5.8$,
$\hat F^{-1}(0.75) = 6.4$.

\begin{align*}
  &w = (7.9 - 4.3)/4 = 3.6/4 = 0.9 && \text{등간격 폭}\\
  &4.3 + 0.9 = 5.2 && \text{첫째 경계}\\
  &4.3 + 1.8 = 6.1 && \text{둘째 경계}\\
  &4.3 + 2.7 = 7.0 && \text{셋째 경계}
\end{align*}

\Cref{tab:qq-iris-bins}가 두 방법을 비교한다.

\begin{table}[ht]\centering
\begin{tabularx}{\linewidth}{@{}l X X X@{}}
\toprule
방법 & 경계 & 폭 & 개수 \\
\midrule
등빈도(오른쪽 닫힘) & 5.1, 5.8, 6.4 & 0.8, 0.7, 0.6, 1.5 & 41, 39, 35, 35 \\
등간격(왼쪽 닫힘)   & 5.2, 6.1, 7.0 & 0.9, 0.9, 0.9, 0.9 & 41, 48, 48, 13 \\
등간격(오른쪽 닫힘) & 5.2, 6.1, 7.0 & 0.9, 0.9, 0.9, 0.9 & 45, 50, 43, 12 \\
\bottomrule
\end{tabularx}
\caption{붓꽃 꽃받침 길이 4구간: 등빈도 대 등간격(붓꽃 150송이).}
\label{tab:qq-iris-bins}
\end{table}

두 교훈. 등간격은 폭이 같고 개수가 다르며(꼬리 구간이 12 또는 13개로
가장 적다), 등빈도는 개수가 비슷하지만(목표 $150/4=37.5$) 폭이 다르다
(0.6 대 1.5). 동점이 많아 등빈도조차 정확히 같은 개수를 주지 않는다.
\end{example}

\begin{pitfall}[두 가지 ``분위수'']
슬라이드의 $\hat F^{-1}(q) = \min\{x \mid \hat F(x) \ge q\}$는 보간 없는
계단 규칙(\texttt{numpy}의 \texttt{method="inverted\_cdf"})이고,
\Cref{ex:quantile}의 규칙은 선형 보간이다. 붓꽃 데이터에서는 두 규칙이
같은 값 5.1, 5.8, 6.4를 주는데, 이는 값이 0.1 격자 위에 있고 그 경계
자리에 동점이 몰려 있기 때문이다(예: $\hat F(5.0)=0.2133$, $\hat
F(5.1)=0.2733$이 $0.25$를 사이에 두고, 선형 보간도 같은 인덱스 구간에서
같은 값을 준다). 다른 데이터에서는 두 규칙이 다른 값을 줄 수 있으므로
어느 규칙을 썼는지 밝힌다.
\end{pitfall}

\begin{pitfall}[등간격 경계가 데이터 값과 겹친다]
붓꽃의 등간격 경계 5.2, 6.1, 7.0은 모두 데이터에 실제로 나오는 값이다.
그 값을 어느 구간에 넣느냐(왼쪽 닫힘 $[\cdot, v_i)$ 대 오른쪽 닫힘
$(\cdot, v_i]$)에 따라 개수가 41, 48, 48, 13 대 45, 50, 43, 12로
갈린다. 경계가 데이터 격자와 우연히 겹치면 이 차이가 커지므로, 이산화
결과를 보고할 때는 닫힌 쪽을 반드시 적는다.
\end{pitfall}

% =====================================================================
\section{범위 정규화와 z-정규화}
\label{sec:normalization}
% source: deck 7 slide 16, 17, 18

\begin{inbrief}
  \item[Why] 단위가 큰 변수가 거리 계산을 독차지한다(소득이 나이를
             압도).
  \item[When] 척도가 크게 다른 변수를 거리 기반 알고리즘에 넣거나 한
              그림에 놓을 때.
  \item[Where] 나이와 소득 10명(슬라이드 18).
  \item[How] 범위 정규화 $x' = (x-\min)/\hat r$는 $[0,1]$로, z-정규화
             $x' = (x-\hat\mu)/\hat\sigma$는 평균 0, 표준편차 1로 옮긴다.
\end{inbrief}

\paragraph{고치는 문제.} 나이(범위 12--40)와 소득(범위 300--6000)을 같은
유클리드 거리 식에 넣으면, 값의 크기가 큰 소득의 차이가 거리를 거의
전부 결정한다. 여러 실험을 합칠 때도 실험마다 기술적 차이(스케일)가
있으면 같은 문제가 생긴다.

\paragraph{아이디어.} 모든 변수를 같은 스케일로 옮긴다. 범위 정규화는
최솟값을 빼고 범위로 나눠 $[0,1]$에 맞추고, z-정규화는 평균을 빼고
표준편차로 나눠 평균 0, 표준편차 1에 맞춘다. 둘 다 더하고 곱하는
\textbf{1차 변환}이라, 값의 상대적 순서와 분포의 모양은 그대로 둔다
(\Cref{sec:qq-apply}의 핵심 문장).

\begin{example}[나이와 소득 10명]
\label{ex:qq-age-income}
슬라이드 18의 값이다. 나이: 12, 14, 18, 23, 27, 28, 34, 37, 39, 40.
소득: 300, 500, 1000, 2000, 3500, 4000, 4300, 6000, 2500, 2700. 둘째
사람($x_2$: 나이 14, 소득 500)의 정규화 값을 본다.

\begin{align*}
  &(40 - 12) = 28 && \text{나이 범위}\\
  &(14 - 12)/28 = 0.0714 && \text{나이 범위 정규화}\\
  &(6000 - 300) = 5700 && \text{소득 범위(최댓값은 } x_8\text{, 6000)}\\
  &(500 - 300)/5700 = 0.0351 && \text{소득 범위 정규화}
\end{align*}
\end{example}

\begin{pitfall}[슬라이드 값은 2400, 다시 계산하면 5700]
슬라이드 18의 식 이미지는 소득 범위를 ``$2700-300=2400$''이라 적었다.
소득의 마지막 행이 2700이라 그 값을 최댓값으로 읽은 것으로 보이지만,
소득의 \textbf{실제 최댓값}은 여덟째 사람의 6000이다. 다시 계산하면
$6000 - 300 = 5700$이다. 화면에 보이는 슬라이드는 강사가 값을 6000,
5700, 5700, 0.035로 이미 덧씌워 고쳤고, $(500-300)/5700 = 0.0351$은
이 5700에서 나오므로 맞다. 식 이미지 자체는 $2400$과 결과 $0.035$가
스스로 맞지 않는다($200/2400 = 0.0833 \ne 0.035$). 교훈은 범위를 잴 때
\textbf{표의 마지막 행이 아니라 실제 최댓값}을 찾는 것이다.
\end{pitfall}

\begin{align*}
  &\bar x_{\text{나이}} = (12{+}14{+}18{+}23{+}27{+}28{+}34{+}37{+}39{+}40)/10 = 27.2 && \text{나이 평균}\\
  &\bar x_{\text{소득}} = 26800/10 = 2680 && \text{소득 평균}\\
  &\hat\sigma_{\text{나이}}(n) = \sqrt{953.6/10} = 9.765 && \text{나이 } \sigma(n)\\
  &\hat\sigma_{\text{나이}}(n{-}1) = \sqrt{953.6/9} = 10.293 && \text{나이 } \sigma(n{-}1)\\
  &\hat\sigma_{\text{소득}}(n) = \sqrt{29{,}796{,}000/10} = 1726.15 && \text{소득 } \sigma(n)\\
  &\hat\sigma_{\text{소득}}(n{-}1) = \sqrt{29{,}796{,}000/9} = 1819.52 && \text{소득 } \sigma(n{-}1)\\
  &z_{\text{나이}}(n) = (14-27.2)/9.765 = -1.352 && \text{나이 } z(n)\\
  &z_{\text{소득}}(n) = (500-2680)/1726.15 = -1.263 && \text{소득 } z(n)\\
  &z_{\text{나이}}(n{-}1) = -13.2/10.293 = -1.282 && \text{나이 } z(n{-}1)\\
  &z_{\text{소득}}(n{-}1) = -2180/1819.52 = -1.198 && \text{소득 } z(n{-}1)
\end{align*}

\begin{pitfall}[슬라이드 값은 $z = (-1.35, -1.26)$, 다시 계산하면
$z = (-1.28, -1.20)$]
슬라이드 17은 $\hat\sigma$를 ``sample variance''라 부르지만 실제 계산은
$n$으로 나눈다(위 $\hat\sigma(n)$). 이 책의 표준편차(\Cref{def:spread})는
$n-1$로 나눈다. 슬라이드 값 $z = (-1.35, -1.26)$은 $n$으로 나눈 것이고,
$n-1$로 다시 계산하면 $z = (-1.28, -1.20)$이다. \texttt{pandas}의
\texttt{.std()} 기본값은 $n-1$, \texttt{numpy}의 \texttt{np.std}
기본값은 $n$이므로, 어느 라이브러리를 쓰든 기본값을 확인해야 한다.
\end{pitfall}

\Cref{tab:qq-scaling}가 $x_2$의 모든 변환 값을 한곳에 모은다.

\begin{table}[ht]\centering
\begin{tabularx}{\linewidth}{@{}l X X@{}}
\toprule
방법 & 나이($x_2$) & 소득($x_2$) \\
\midrule
원래 값 & 14 & 500 \\
범위 정규화 & 0.071 & 0.035 \\
z-정규화($n$, 슬라이드) & $-1.35$ & $-1.26$ \\
z-정규화($n-1$, 책) & $-1.28$ & $-1.20$ \\
로버스트(\Cref{ex:qq-robust}) & $-0.794$ & $-0.800$ \\
\bottomrule
\end{tabularx}
\caption{나이와 소득 $x_2$의 변환 값 한눈에 보기.}
\label{tab:qq-scaling}
\end{table}

\Cref{fig:qq-scaling}가 원래, 범위, z($n-1$), 로버스트 네 패널을
산점도로 보여준다.

\begin{figure}[ht]\centering
  \includegraphics[width=0.95\linewidth]{07_07_age_income_scaling.png}
  \caption{나이(가로)와 소득(세로) 10명의 산점도 네 패널. 원래 패널은
  축 비율을 그대로 둬 소득 축이 나이 축을 압도함을 보인다. $x_2$를
  강조색으로 표시했다.}
  \label{fig:qq-scaling}
\end{figure}

% =====================================================================
\section{분위수 정규화}
\label{sec:quantile-norm}
% source: deck 7 slide 16, 19, 20

\begin{inbrief}
  \item[Why] 실험(표본)마다 측정 장비의 기술적 차이로 분포 \textbf{전체}가
             어긋나면, 평균과 표준편차만 맞춰서는 부족하다.
  \item[When] 여러 표본이 같은 분포를 가져야 한다고 믿을 근거가 있을 때
              (유전자 발현 배열).
  \item[Where] 유전자 4개 $\times$ 표본 4개 행렬.
  \item[How] 열마다 정렬 $\to$ 같은 순위끼리 평균 $\to$ 그 평균을 원래
             순위 자리에 되돌린다.
\end{inbrief}

\paragraph{고치는 문제.} z-정규화(\Cref{sec:normalization})는 평균과
표준편차만 맞춘다. 그러나 두 표본의 분포 \textbf{모양} 자체가 다르면
(한쪽은 치우치고 한쪽은 대칭이면), 평균과 표준편차를 맞춰도 분위수마다
어긋남이 남는다. 실험 간 기술적 차이(배치 효과, Batch Effect)가 이런
모양 차이를 만든다.

\paragraph{아이디어.} 모든 표본(열)의 분위수를 통째로 같게 만든다. 열마다
정렬한 뒤, 같은 순위 자리(가장 작은 값끼리, 둘째로 작은 값끼리, ...)의
평균을 구하고, 그 평균값을 원래 순위 자리에 되돌린다. 그러면 모든 열의
값 집합(분위수)이 똑같아지고, 열 \textbf{안의} 순서(어느 유전자가 더
큰가)는 보존된다.

\begin{example}[분위수 정규화 손 계산]
\label{ex:qq-quantile-norm}
슬라이드 20의 $4 \times 4$ 행렬이다(행 = 유전자, 열 = 표본).

\[
  \begin{pmatrix} 1 & 5 & 3 & 5 \\ 2 & 1 & 6 & 7 \\ 3 & 2 & 2 & 6 \\ 4 & 6 & 1 & 8 \end{pmatrix}
\]

\begin{align*}
  &\text{열 정렬: } (1,2,3,4),\ (1,2,5,6),\ (1,2,3,6),\ (5,6,7,8) && \text{네 열}\\
  &(1{+}1{+}1{+}5)/4 = 2 && \text{1위 평균}\\
  &(2{+}2{+}2{+}6)/4 = 3 && \text{2위 평균}\\
  &(3{+}5{+}3{+}7)/4 = 4.5 && \text{3위 평균}\\
  &(4{+}6{+}6{+}8)/4 = 6 && \text{4위 평균}
\end{align*}

\Cref{tab:qq-quantile-norm}이 원래, 정렬, 순위별 평균, 되돌린 결과 네
단계를 보여준다.

\begin{table}[ht]\centering
\begin{tabularx}{\linewidth}{@{}l X X X X@{}}
\toprule
단계 & 열1 & 열2 & 열3 & 열4 \\
\midrule
원래 & 1, 2, 3, 4 & 5, 1, 2, 6 & 3, 6, 2, 1 & 5, 7, 6, 8 \\
정렬 & 1, 2, 3, 4 & 1, 2, 5, 6 & 1, 2, 3, 6 & 5, 6, 7, 8 \\
순위별 평균 & 2, 3, 4.5, 6 & 2, 3, 4.5, 6 & 2, 3, 4.5, 6 & 2, 3, 4.5, 6 \\
되돌린 결과 & 2, 3, 4.5, 6 & 4.5, 2, 3, 6 & 4.5, 6, 3, 2 & 2, 4.5, 3, 6 \\
\bottomrule
\end{tabularx}
\caption{$4\times4$ 분위수 정규화 손 계산(슬라이드 20). 16칸 모두
슬라이드와 일치한다.}
\label{tab:qq-quantile-norm}
\end{table}

결과 행렬의 네 열 모두 값 집합 $\{2, 3, 4.5, 6\}$을 갖는다. 열1은 원래
값이 이미 순위 순서였으므로 정렬 결과와 되돌린 결과가 같다.
\end{example}

\begin{pitfall}[동점이 있으면 순위 규칙이 필요하다]
슬라이드 20의 행렬에는 동점이 없어 되돌리기가 깔끔하다. 동점이 있으면
어떻게 나눌지 규칙이 필요하다. 표본 두 개, 유전자 세 개로 확인한다.
$A = (5, 2, 3)$, $B = (4, 4, 1)$(책의 예시 값). 정렬 $A: 2, 3, 5$,
정렬 $B: 1, 4, 4$. 순위별 평균은 $(2{+}1)/2 = 1.5$, $(3{+}4)/2 = 3.5$,
$(5{+}4)/2 = 4.5$다. $A$를 되돌리면 $5$는 3위, $2$는 1위, $3$은 2위이므로
$(4.5, 1.5, 3.5)$다. $B$의 $1$은 1위라 $1.5$를 받지만, 두 개의 $4$는
2위와 3위를 나눠 가진다. (i) 나온 순서대로 주면 $(3.5, 4.5, 1.5)$가
되어 \textbf{같은 값 4가 서로 다른 값이 되는} 문제가 생긴다. (ii)
두 자리의 평균 $(3.5+4.5)/2 = 4.0$을 둘 다에 주면 $(4.0, 4.0, 1.5)$로
동점이 보존된다. 이 책은 (ii)를 권하고, 어느 규칙을 썼는지 밝혀 쓰라고
적는다.
\end{pitfall}

\Cref{fig:qq-quantile-norm}이 원래 행렬과 정규화 결과를 나란히 그린다.

\begin{figure}[ht]\centering
  \includegraphics[width=0.9\linewidth]{07_08_quantile_norm.png}
  \caption{왼쪽: 원래 $4\times4$ 행렬, 표본(열)별 점. 오른쪽: 분위수
  정규화 뒤, 네 표본의 값 집합이 $\{2, 3, 4.5, 6\}$으로 같아졌다.
  같은 유전자(행)는 같은 색으로 이었다.}
  \label{fig:qq-quantile-norm}
\end{figure}

% =====================================================================
\section{로버스트 스케일링}
\label{sec:robust-scaling}
% source: deck 7 slide 21

\begin{inbrief}
  \item[Why] 평균과 표준편차는 극단값 하나에 끌려간다(\Cref{def:spread}).
  \item[When] 극단값이나 입력 실수가 의심될 때.
  \item[Where] 소득에 잘못 입력된 60000(책의 가정 예).
  \item[How] $x' = (x - \text{중앙값})/\IQR$.
\end{inbrief}

\paragraph{고치는 문제.} z-정규화의 $\hat\mu, \hat\sigma$는 모든 값을
평균에 반영한다. 값 하나가 극단적으로 크거나 작으면 평균과 표준편차
모두 그 값 쪽으로 끌려가고, 정규화 결과 전체가 비뚤어진다.

\paragraph{아이디어.} 평균 대신 \textbf{중앙값}, 표준편차 대신
\textbf{사분위 범위(IQR)}를 쓴다. 둘 다 가운데 순위에서 정해지므로
극단값 하나가 아무리 커져도 값 자체는 거의 움직이지 않는다(\term{로버스트}
(Robust)의 뜻).

\begin{example}[로버스트 스케일링과 입력 실수]
\label{ex:qq-robust}
사분위수는 \Cref{ex:quantile}의 선형 보간 규칙을 쓴다.

나이 정렬: 12, 14, 18, 23, 27, 28, 34, 37, 39, 40.

\begin{align*}
  &\text{중앙값} = (27+28)/2 = 27.5 && \text{나이}\\
  &Q_1: \text{인덱스 } 2.25 \to 18 + 0.25(23{-}18) = 19.25 && \text{나이 } Q_1\\
  &Q_3: \text{인덱스 } 6.75 \to 34 + 0.75(37{-}34) = 36.25 && \text{나이 } Q_3\\
  &\IQR = 36.25 - 19.25 = 17 && \text{나이 } \IQR\\
  &x_2' = (14 - 27.5)/17 = -0.794 && \text{나이 } x_2\text{의 로버스트 값}
\end{align*}

소득 정렬: 300, 500, 1000, 2000, 2500, 2700, 3500, 4000, 4300, 6000.

\begin{align*}
  &\text{중앙값} = (2500+2700)/2 = 2600 && \text{소득}\\
  &Q_1 = 1000 + 0.25(2000{-}1000) = 1250 && \text{소득 } Q_1\\
  &Q_3 = 3500 + 0.75(4000{-}3500) = 3875 && \text{소득 } Q_3\\
  &\IQR = 3875 - 1250 = 2625 && \text{소득 } \IQR\\
  &x_2' = (500-2600)/2625 = -0.800 && \text{소득 } x_2\text{의 로버스트 값}
\end{align*}

이제 여덟째 사람의 소득 6000을 60000으로 잘못 입력했다고 하자(책의
가정). 정렬 위치 7까지의 값(중앙값, $Q_1$, $Q_3$)은 바뀌지 않으므로

\begin{align*}
  &\text{로버스트 값} = -0.800 && \text{그대로(순위가 안 바뀜)}\\
  &\bar x_{\text{새}} = (26800 - 6000 + 60000)/10 = 8080 && \text{z-정규화 새 평균}\\
  &\hat\sigma_{\text{새}}(n) = \sqrt{301{,}275{,}600} = 17357.3 && \text{z-정규화 새 } \sigma\\
  &x_2' = (500-8080)/17357.3 = -0.437 && \text{z, 원래} -1.263\text{에서 이동}\\
  &(500-300)/(60000-300) = 200/59700 = 0.00335 && \text{범위, 원래 } 0.035\text{에서} \tfrac1{10}\text{로}
\end{align*}

입력 실수 하나가 z-정규화 값은 $-1.26$에서 $-0.44$로, 범위 정규화 값은
10분의 1로 바꾸지만, 로버스트 값은 전혀 움직이지 않는다.
\end{example}

% =====================================================================
\section{로그 변환}
\label{sec:log-transform}
% source: deck 7 slide 22

\begin{inbrief}
  \item[Why] 오른쪽으로 크게 치우친 변수는 패턴이 한쪽에 뭉쳐 보이고
             정규성 가정을 어긴다.
  \item[When] 값이 모두 양수이고 여러 자릿수에 걸쳐 퍼질 때.
  \item[Where] 나라별 인구와 면적(슬라이드 그림), 당류(\texttt{sugars}).
  \item[How] $\log x$를 쓰고, 0이 있으면 $\log(1+x)$를 쓴다.
\end{inbrief}

\paragraph{고치는 문제.} 값이 여러 자릿수에 걸쳐 퍼진 변수(인구, 소득,
당류)는 큰 값 몇 개가 그림과 평균을 지배하고, 작은 값들은 한곳에
뭉쳐 보인다. 오른쪽 치우침(\Cref{sec:skew-sign})이 크면 정규성을
가정하는 방법이 잘 맞지 않는다.

\paragraph{아이디어.} 로그는 큰 값의 간격을 압축하고 작은 값의 간격을
상대적으로 펼친다($\log 1000 - \log 100 = \log 100 - \log 10$, 같은
\textbf{배율} 차이는 로그축에서 같은 거리다). 오른쪽 꼬리를 눌러
왜도를 줄이는 효과가 있다.

\begin{pitfall}[$\log 0$은 정의되지 않는다]
패스트푸드 \texttt{sugars}는 17개 항목이 0이다(\Cref{ex:zero-as-missing}).
이 값들은 실제로 당류가 없는 음식이고 결측이 아니다. $\log x$를 그대로
쓰면 이 17행이 $-\infty$가 되거나 소프트웨어가 조용히 결측으로 바꿔
버린다. 대신 $\log(1+x)$(\texttt{numpy.log1p})를 쓴다. 0을 결측으로
바꿔 ``해결''하는 것은 \Cref{ex:zero-as-missing}의 오류를 되풀이하는
것이다.
\end{pitfall}

\begin{align*}
  &\text{왜도}(\text{sugars}) = 2.7088 && \text{로그 전(오른쪽으로 크게 치우침)}\\
  &\text{왜도}(\log(1+\text{sugars})) = 0.0146 && \text{로그 후}
\end{align*}

\begin{pitfall}[왜도가 0에 가깝다고 정규가 된 것은 아니다]
$\log(1+\text{sugars})$의 왜도는 $0.0146$으로 0에 매우 가깝지만, 이
분포는 여전히 \textbf{여러 봉우리}다. 0이던 17개 항목이 $\log(1+0)=0$
한 지점에 모이고, 나머지 값들이 4.3 근처에 또 뭉친다. 왜도 하나만
보고 ``로그 변환으로 정규분포가 됐다''고 쓰면 안 된다. 왜도는 \textbf{좌우
대칭}만 재고, 봉우리 개수는 재지 않는다.
\end{pitfall}

\Cref{fig:qq-log}가 변환 전후의 히스토그램과 Q-Q를 함께 보여준다.

\begin{figure}[ht]\centering
  \includegraphics[width=0.95\linewidth]{07_09_sugars_log.png}
  \caption{\texttt{sugars} 히스토그램과 정규 Q-Q(위), $\log(1+x)$ 변환
  후(아래). 왜도는 $2.7088 \to 0.0146$으로 줄지만, 0인 17개 항목이
  만드는 봉우리 때문에 Q-Q는 여전히 휜다.}
  \label{fig:qq-log}
\end{figure}

% =====================================================================
\section{결측값 처리}
\label{sec:impute}
% source: deck 7 slide 24, 25 (+ deck 2 slide 9-11)

\begin{inbrief}
  \item[Why] 빈칸이 있으면 대부분의 계산이 실패하거나 조용히 행을
             버린다.
  \item[When] \texttt{isna()}로 빈칸을 찾고, 그 빈칸이 진짜 결측인지
              (0이 아닌지) 확인한 뒤.
  \item[Where] 수치형(평균 대체), 범주형(최빈값 대체).
  \item[How] 행이나 열을 삭제하거나, 평균, 최빈값, 자리표시자
             (Placeholder), 예측값으로 채운다.
\end{inbrief}

채우기 전에 그 빈칸이 진짜 결측인지부터 확인한다(\Cref{ex:zero-as-missing}).
0을 결측으로 오해하면, 채우는 방법을 아무리 잘 골라도 시작부터 틀린다.
같은 원리가 슬라이드에도 나온다. 덱 2 슬라이드 9는 완전성(Completeness,
주제를 빠짐없이 담았는가)을 정제된 데이터(Clean Data)의 기준 하나로
든다. 슬라이드 10(``Zeros replace missing values'')은 그 기준이 깨지는
반대 방향의 실수, 즉 진짜 결측이 0으로 채워져 있는 경우를 경고한다.
슬라이드 11은 같은 계열의 다른 결함(범주가 나쁘게 정의됨, 값이 의심스러움,
데이터가 너무 거칢)을 나열해, 결측은 데이터 정제에서 마주치는 여러
문제 중 하나임을 보여준다.

\paragraph{고치는 문제.} 결측값이 있는 행은 평균, 표준편차 같은 대부분의
계산에서 오류를 내거나 \texttt{NaN}을 전파한다. 그 행을 그냥 지우면
표본 크기 $n$이 줄고, 데이터를 얻는 데 든 비용(슬라이드 25의 ``less
training data'')이 사라진다.

\paragraph{아이디어.} 선택은 셋이다. \textbf{삭제}(행 또는 열을 지움,
$n$을 잃음), \textbf{대체}(평균, 최빈값, 자리표시자로 채움, $n$은
지키지만 퍼짐을 줄임), \textbf{예측}(다른 변수로 값을 추정, 이 덱에는
방법이 없음). 무엇을 고르든 대가가 있다.

\begin{example}[평균 대체의 대가]
\label{ex:qq-impute}
여덟째 사람의 소득 6000이 비었다고 하자(책의 가정, \Cref{ex:qq-age-income}
자료).

\begin{align*}
  &\text{남은 9개 합} = 26800 - 6000 = 20800 && \text{소득 합에서 뺌}\\
  &\text{평균}(n{=}9) = 20800/9 = 2311.1 && \text{빈칸 뺀 평균}\\
  &\text{중앙값}(n{=}9) = 2500 && \text{정렬 5번째(300..4300)}\\
  &\text{중앙값}(n{=}10,\text{평균 채움}) = (2311.1{+}2500)/2 = 2405.6 && \text{정렬 5,6번째 평균}\\
  &\text{sd 진짜}(n{=}10) = 1819.5 && \text{원래 10개}\\
  &\text{sd 빈칸 뺀}(n{=}9) = \sqrt{17{,}548{,}889/8} = 1481.1 && \text{9개로 계산}\\
  &\text{sd 평균 채움}(n{=}10) = \sqrt{17{,}548{,}889/9} = 1396.4 && \text{채운 값은 편차 0}
\end{align*}

채운 값(평균)은 평균과의 거리가 0이므로 제곱편차 합은 9개일 때와 같은데,
나누는 수만 $n-1$ 기준으로 8에서 9로 커져 표준편차가 더 작아진다.

\Cref{tab:qq-impute}가 방법별로 정리한다.

\begin{table}[ht]\centering
\begin{tabularx}{\linewidth}{@{}l X X X@{}}
\toprule
방법 & $n$ & 평균/중앙값 & 표준편차 \\
\midrule
진짜(빈칸 없음) & 10 & 2680 / 2600 & 1819.5 \\
삭제 & 9 & 2311.1 / 2500 & 1481.1 \\
평균 대체 & 10 & 2311.1(채운 값) / 2405.6 & 1396.4 \\
\bottomrule
\end{tabularx}
\caption{소득 결측 처리 방법별 $n$, 평균, 표준편차 비교.}
\label{tab:qq-impute}
\end{table}
\end{example}

\begin{pitfall}[평균 대체는 표본 크기를 지키지만 퍼짐을 줄인다]
평균 대체는 $n=10$을 지킨다는 점에서 삭제보다 나아 보이지만, 표준편차가
$1819.5 \to 1396.4$로 줄고, 가장 큰 값(6000)이 있었다는 사실 자체가
지워진다. 삭제는 $n$을 잃지만($10 \to 9$) 남은 값의 표준편차($1481.1$)는
평균 대체보다 진짜 값에 더 가깝다. 어느 쪽도 공짜가 아니다.
\end{pitfall}

% =====================================================================
\section{어떻게 적용하는가}
\label{sec:qq-apply}
% source: deck 7 slide 23

\Cref{tab:qq-transforms}가 문제와 변환을 짝짓는다.

\begin{table}[ht]\centering
\begin{tabularx}{\linewidth}{@{}l X X X@{}}
\toprule
문제 & 쓰는 변환 & 바뀌는 것 & 바뀌지 않는 것 \\
\midrule
척도 차이 & 범위/z-정규화 & 위치와 폭 & 모양 \\
극단값 & 로버스트 스케일링 & 위치와 폭 & 모양 \\
긴 오른쪽 꼬리 & 로그 변환 & 모양 & 순서 \\
표본 간 분포 차이 & 분위수 정규화 & 분포 전체 & 열 안의 순서 \\
범주가 필요 & 이산화 & 값이 등급으로 & 등급 사이 순서 \\
\bottomrule
\end{tabularx}
\caption{문제에 따라 쓰는 변환과 그 효과.}
\label{tab:qq-transforms}
\end{table}

\begin{keyidea}[1차 변환은 모양을 바꾸지 않는다]
범위 정규화, z-정규화, 로버스트 스케일링은 모두 더하고 곱하는 \textbf{1차
변환}이다. Q-Q 플롯의 점은 직선 위든 아니든 \textbf{그대로}이고, 축
눈금만 바뀐다. 모양(치우침, 꼬리 두께)을 실제로 바꾸는 것은 로그
변환과 분위수 정규화뿐이다.
\end{keyidea}

\Cref{fig:qq-flow}가 이 장의 절차를 한 장으로 정리한다.

\begin{figure}[ht]\centering
\resizebox{\linewidth}{!}{%
\begin{tikzpicture}[
  x = 1mm, y = 1mm,
  flowbox/.style = {draw = bookrule, fill = booksurface, rounded corners = 2pt,
                   minimum width = 34mm, minimum height = 10mm, align = center,
                   font = \small, inner sep = 1pt},
  leafbox/.style = {draw = bookaccent, fill = bookaccentbg, rounded corners = 2pt,
                   minimum width = 30mm, minimum height = 10mm, align = center,
                   font = \small, inner sep = 1pt},
  arrow/.style  = {-{Stealth[length=2.2mm]}, draw = bookmuted}]
  \node[flowbox] (hist) at (0, 0)   {히스토그램 + 후보};
  \node[flowbox] (cdf)  at (46, 0)  {CDF 겹치기};
  \node[flowbox] (qq)   at (92, 0)  {Q-Q 플롯};
  \node[flowbox] (need) at (92, -20) {변환이 필요한가?};
  \node[leafbox] (no)   at (46, -34) {아니오\\[-1pt]\footnotesize 분석 진행};
  \node[flowbox] (pick) at (138, -20) {\Cref{tab:qq-transforms}로\\[-1pt]변환 고르기};
  \node[flowbox] (redo) at (138, -40) {변환 뒤\\[-1pt]Q-Q 다시 보기};
  \draw[arrow] (hist) -- (cdf);
  \draw[arrow] (cdf) -- (qq);
  \draw[arrow] (qq) -- (need);
  \draw[arrow] (need) -- node[above, sloped, font=\footnotesize]{아니오} (no);
  \draw[arrow] (need) -- node[above, sloped, font=\footnotesize]{예} (pick);
  \draw[arrow] (pick) -- (redo);
\end{tikzpicture}%
}
\caption{이 장의 절차. 확인(히스토그램 $\to$ CDF $\to$ Q-Q) 뒤 변환
필요 여부를 판단하고, 변환했으면 Q-Q로 다시 확인한다.}
\label{fig:qq-flow}
\end{figure}

이 책 안에서 이 장의 절차가 쓰이는 자리는 다음과 같다.

\begin{itemize}
  \item \Cref{fig:normal-fit}와 \Cref{fig:qq-calories}: 패스트푸드
        \texttt{calories}의 히스토그램은 종 모양처럼 보이지만, Q-Q
        플롯과 Shapiro-Wilk 검정은 정규성을 기각한다($p = 0.00038$).
  \item \texttt{sugars}: 로그 변환 전후 왜도($2.7088 \to 0.0146$)와
        Q-Q(\Cref{fig:qq-log}, \Cref{sec:skew-sign}).
  \item \Cref{sec:hw2-compare-how}, \Cref{fig:hw2-movie-normal}: HW2가
        이미 QQ 도표를 썼다. \Cref{tab:qq-shapes}로 다시
        읽으면 ``양 끝만 어긋난다''는 왼쪽 꼬리가 더 긴 약한 왼쪽 치우침
        (왜도 $-0.48$)이라고 이름 붙일 수 있다. 새 계산은 하지 않는다.
  \item \Cref{sec:cd-t-one}, \Cref{sec:cd-anova}: t검정과 F검정의
        정규성 가정을 Q-Q로 먼저 확인하는 순서.
  \item \Cref{sec:cd-ks}, \Cref{sec:hw2-learned}: Q-Q는 \textbf{어디서}
        어긋나는지, KS $D$는 \textbf{얼마나} 어긋나는지(숫자 하나)를
        보여준다. 둘을 함께 쓴다.
  \item \Cref{ch:project}: 분석 전 점검 목록(분포 확인 $\to$ 변환
        $\to$ 결측)으로 이 장의 순서를 그대로 쓴다.
\end{itemize}

\Cref{fig:qq-calories}가 \texttt{calories}의 Q-Q와 Shapiro-Wilk p값을
함께 보여준다.

\begin{figure}[ht]\centering
  \includegraphics[width=0.7\linewidth]{07_10_calories_qq.png}
  \caption{패스트푸드 \texttt{calories}의 정규 Q-Q 플롯. Shapiro-Wilk
  $p = 0.00038 < 0.05$로 정규성을 기각한다.}
  \label{fig:qq-calories}
\end{figure}

확인 문제 셋(슬라이드 23의 Kahoot 대신).

\begin{enumerate}
  \item 점이 양 끝에서 선 위로 휘어 오르는 Q-Q는 어느 모양인가?
        (답: 오른쪽 치우침, \Cref{tab:qq-shapes})
  \item 붓꽃 꽃받침 길이를 $k=4$ 등간격으로 나누면 경계는?
        (답: 5.2, 6.1, 7.0, \Cref{ex:qq-iris})
  \item 소득 500을 범위 정규화하면 값은?
        (답: 0.035, \Cref{ex:qq-age-income})
\end{enumerate}

\subsection{붓꽃 데이터 출처}
\label{sec:qq-iris-source}

슬라이드 15는 데이터 이름을 ``sepal length''로만 적고 출처를 밝히지
않는다. $n=150$(개수 합), 범위 $[4.3, 7.9]$, 0.1 단위 값은 피셔의
붓꽃 데이터(Fisher's Iris Data, 3종 각 50송이)의 꽃받침 길이와 맞는다.
이 책은 \texttt{sklearn.datasets.load\_iris()}로 값만 읽는다(모형 학습
없음). 슬라이드 14의 점들은 등간격·등빈도 개념을 보여주는 장난감
데이터이고 붓꽃이 \textbf{아니다}.

% =====================================================================
\section{앞으로}
\label{sec:qq-future}
% source: deck 7 slide 26

\begin{itemize}
  \item \textbf{정규성 전용 검정과 모수 추정 뒤의 KS 검정.} Shapiro-Wilk
        같은 정규성 전용 검정이 왜 KS 검정과 다른 답을 줄 수 있는지는
        이 덱에 없다. \Cref{sec:cd-future}가 같은 질문을 열어 두었다.
  \item \textbf{분석 방법 선택 자체가 결론을 바꾼다.} 같은 축구 경기
        데이터를 받은 29개 연구 팀 중 20팀은 유의, 9팀은 유의하지
        않음을 보고했다(그림에서 센 값). 그림에 적힌 출처는
        FiveThirtyEight, Brian Nosek 등이다(슬라이드 26 그림 안의
        인쇄된 문구를 그대로 옮긴다, 새 참고문헌 항목은 만들지 않는다).
        \Cref{sec:pv-threshold}와 연결된다. 분석 선택이 많을수록 우연한
        유의가 늘어난다.
  \item \textbf{고급 대체(Advanced Imputation).} 예측 모형으로 결측을
        채우는 방법(슬라이드 25). 이 덱에는 방법이 없으므로 열린
        질문으로 남긴다.
  \item \textbf{분위수 정규화의 위험.} ``분포가 같아야 한다''는 전제가
        틀리면, 분위수 정규화가 진짜 생물학적 차이까지 지울 수 있다.
        이 덱에는 다루지 않으므로 열린 질문으로 남긴다.
\end{itemize}
